\begin{abstract}
LLM trustworthiness dimensions exhibit sparse coupling---failures correlate in only 2 of 10 dimension pairs, not pervasively. We demonstrate a methodology for detecting and characterizing coupling patterns via phi coefficient analysis across 5 dimensions (truthfulness, robustness, fairness, safety, privacy) on synthetic benchmark emulation (n=1500 instances across 3 simulated model profiles: GPT-4, Claude-3, Llama-3). Multi-dimensional evaluation frameworks assume dimensional independence, yet this masks compound failures---instances failing multiple criteria simultaneously. Our synthetic validation shows coupling is detectable when present (phi 0.33--0.40, $p < 1\times10^{-13}$) but sparse: zero models exhibit $\geq$3 coupled pairs after multiple comparison correction. Partial correlation controlling for instance difficulty yields retention 86--161\%, establishing that coupling reflects shared vulnerability mechanisms rather than spurious artifacts of hard instances failing everywhere. Models show suggestive but statistically inconclusive coupling profile differences (Mantel $r < 0.7$, $p > 0.0167$), requiring larger samples ($n \geq 500$) for confirmation. \textbf{Synthetic data limitation:} All findings validate measurement methodology; real-world coupling characterization pending production benchmark access. If validated, implications include: benchmarks should report coupling statistics alongside per-dimension scores, and model selection can exploit known patterns for deployments requiring multiple guarantees.
\end{abstract}
